\chapter{Phenytoin Level}

\section*{What Is This Test?}
The phenytoin level measures the serum concentration of phenytoin, a hydantoin anticonvulsant used for the treatment of generalized tonic-clonic (grand mal) seizures and focal (partial) seizures. Phenytoin stabilizes neuronal membranes by binding to and prolonging the inactivated state of voltage-gated sodium channels, preventing the high-frequency repetitive firing of action potentials that underlies seizure propagation \cite{brodie1996antiepileptic}. Phenytoin exhibits nonlinear (zero-order, Michaelis-Menten) pharmacokinetics at therapeutic concentrations: the hepatic CYP2C9 (and to a lesser extent CYP2C19) enzyme system that metabolizes phenytoin becomes saturated in the therapeutic range, and small dose increases produce disproportionately large increases in serum concentrations \cite{anderson2008pharmacokinetic}. Phenytoin has a narrow therapeutic index (total phenytoin 10--20 mcg/mL). In patients with hypoalbuminemia (nephrotic syndrome, cirrhosis, malnutrition, pregnancy), the free (unbound, pharmacologically active) fraction of phenytoin is increased because phenytoin is approximately 90\% protein-bound, primarily to albumin. In these patients, total phenytoin levels underestimate the pharmacologically active free drug concentration, and a corrected phenytoin level or direct free phenytoin measurement is required.

\section*{Alternative Names}
Dilantin Level; Phenytoin Serum Concentration; Total Phenytoin; Free Phenytoin; Diphenylhydantoin Level

\section*{Principle}
Phenytoin is measured by particle-enhanced turbidimetric inhibition immunoassay (PETINIA) or enzyme-multiplied immunoassay technique (EMIT). The assay uses anti-phenytoin antibodies. Total phenytoin is measured from serum (serum separator tube). Free (unbound) phenytoin is measured after ultrafiltration; the free fraction is normally approximately 10\% (0.1). The corrected phenytoin level for hypoalbuminemia is calculated using the Sheiner-Tozer equation: Corrected phenytoin = Measured total phenytoin / [(0.2 $\times$ albumin) + 0.1] \cite{burton2006applied}.

\section*{History}
Phenytoin was first synthesized in 1908 by the German chemist Heinrich Biltz but was not used clinically until 1938, when Tracy Putnam and Houston Merritt screened over 700 compounds for anticonvulsant activity and identified phenytoin as a potent anticonvulsant without the sedative properties of phenobarbital. The Sheiner-Tozer equation for corrected phenytoin was published in 1976 \cite{sheiner1980clinical}. Fosphenytoin (Cerebyx), a water-soluble phosphate ester prodrug of phenytoin, was introduced in the 1990s to overcome the difficulties of IV phenytoin administration (propyleneglycol vehicle causing hypotension, cardiac arrhythmias, and purple glove syndrome with IV extravasation).

\section*{Why This Test Is Ordered}
Monitoring of phenytoin therapy to ensure levels are within the therapeutic range. Diagnosis of phenytoin toxicity. Dose adjustment. Monitoring during pregnancy (increased clearance and decreased albumin decrease total levels). Monitoring in patients with hypoalbuminemia, renal failure (uremia decreases protein binding), and concurrent protein-binding-displacing drugs (valproate, salicylates, sulfonamides).

\section*{Is This a Primary or Secondary Test}
TDM is a secondary monitoring test. Phenytoin is dosed and titrated on the basis of seizure control and clinical tolerability; the serum level is a secondary tool used to confirm therapeutic dosing, to guide dose changes given its nonlinear (zero-order) pharmacokinetics, and to evaluate suspected toxicity, non-adherence, pregnancy-related changes, or drug interactions.

\section*{How This Test Is Performed}
A blood sample is collected by standard venipuncture into a serum separator tube (SST). Total phenytoin is measured in serum by immunoassay (PETINIA or EMIT) on automated analyzers; free (unbound) phenytoin is measured after ultrafiltration of serum when indicated. Results are typically available within 1--2 hours for urgent (STAT) requests and within 1--2 days for routine samples. After fosphenytoin administration, results are reported in phenytoin equivalents and should be drawn only after the prodrug has fully converted (generally at least 2 hours after the end of an infusion).

\section*{What Preparation Is Needed}
The sample should be drawn as a trough, immediately before the next scheduled dose, at steady state (5--10 days after initiation or dose change), and the daily dosing schedule should be consistent \cite{patsalos2008}. Because phenytoin exhibits nonlinear pharmacokinetics, the exact dose and route (oral vs. IV vs. fosphenytoin), formulation, and time of the last dose should be documented \cite{patsalos2013saliva}. In patients with hypoalbuminemia, renal failure, or pregnancy, a free phenytoin level should be requested, or the total phenytoin corrected with the Sheiner-Tozer equation using a simultaneously drawn albumin.

\section*{Contraindications and Precautions}
\begin{itemize}
  \item No absolute contraindications to standard venipuncture
  \item The result cannot be interpreted without knowing albumin and renal function: hypoalbuminemia, uremia, pregnancy, and displacing drugs (valproate, salicylates, sulfonamides) raise the free fraction, so total levels may appear therapeutic while the free (active) level is toxic
  \item Valproate coadministration both displaces phenytoin from albumin and inhibits its metabolism, so even small added doses can substantially raise free phenytoin
  \item Nonlinear (Michaelis-Menten) kinetics mean small dose increments can cause disproportionate level increases; repeat levels are required at the new steady state after every dose change
  \item After fosphenytoin, samples drawn too soon after IV administration may not reflect complete conversion to phenytoin
\end{itemize}

\section*{Risks}
Minimal risk. Slight pain or bruising at the venipuncture site. Rarely: hematoma, infection, or vasovagal syncope.

\section*{What Happens After the Test}
The result is interpreted against the therapeutic range (total 10--20 mcg/mL; free 1.0--2.0 mcg/mL, or a Sheiner-Tozer corrected value in hypoalbuminemia). Subtherapeutic levels prompt a dose increase with a repeat level after 5--10 days at the new dose; toxic levels (nystagmus, ataxia, lethargy) prompt dose reduction and rechecking until the level is within range. Because metabolism is saturable, small dose changes are recommended and follow-up levels are essential after every adjustment.

\section*{Understanding Results}
Subtherapeutic (<10 mcg/mL): Inadequate seizure control, non-adherence, pregnancy, or drug interactions increasing clearance. Toxic (>20 mcg/mL): Nystagmus (earliest sign, 20--30 mcg/mL), ataxia (30--40 mcg/mL), lethargy and dysarthria (40--50 mcg/mL), coma and seizures (>50 mcg/mL). Long-term toxicity: gingival hyperplasia (50\% of patients), hirsutism, coarsening of facial features, cerebellar atrophy, peripheral neuropathy, megaloblastic anemia (folate deficiency), osteomalacia (vitamin D deficiency).

\section*{Normal Results}
Therapeutic range (total phenytoin): 10--20 mcg/mL. Free phenytoin therapeutic range: 1.0--2.0 mcg/mL. Corrected phenytoin (hypoalbuminemia): calculated using Sheiner-Tozer equation. Timing: Trough level (immediately before next dose) at steady state (5--10 days after initiation or dose change). Fosphenytoin: phenytoin equivalents measured.

\section*{Follow-up Tests}
Serum albumin (for correction of total phenytoin in hypoalbuminemia). Free phenytoin level (if albumin <2.5--3.0 g/dL, renal failure, pregnancy, or clinical toxicity despite normal total phenytoin). Liver function tests (phenytoin hepatotoxicity is rare). Complete blood count (megaloblastic anemia). Vitamin D and bone density (long-term therapy).

\section*{References}
\bibliographystyle{unsrtnat}
\bibliography{references}

\clearpage
